\input{Content/02_Chapter1/Text/diagrams}
\chapter{الإطار المفاهيمي والنظري لعلم اللسانيات الجنائية}

\section*{تمهيد}

قبل أي شيء، نريد أن نضبط توقعات هذا الفصل. ليس الهدف منه حفظ مصطلحات، بل بناء أرضية مفاهيمية تجعل متابعة الفصلين القادمين أيسر وأكثر عمقاً. كما أن إشكالية المصطلح --- هل نقول لسانيات جنائية أم لسانيات قضائية أم لسانيات جنائية قانونية? --- ليست ترفاً فكرياً، بل هي مسألة تمس صميم التواصل العلمي بين الباحثين، وسنعالجها بتفصيل.

ما نحاول فعله هنا هو استعراض ماهية هذا العلم، وتتبع كيف وضع الباحثون حدوده ودرسوا علاقته بالتخصصات المجاورة، ثمّ فهم لماذا يكتسب الدليل اللغوي أهمية متزايدة في منظومة العدالة الحديثة.



%% ======================================================
\clearpage

\section{مفهوم اللسانيات الجنائية}
%% ======================================================

\subsection{تعريف اللسانيات}

\subsubsection*{أولاً: التعريف اللغوي}

يرجع أصل كلمة «اللسانيات» إلى الجذر اللغوي (ل س ن) الذي يدل على اللسان، وهو جارحة الكلام ووسيلة التعبير. وقد يُستعمل اللسان مجازاً للدلالة على الكلام أو اللغة التي يتحدث بها الإنسان أو الجماعة، ومن ذلك قولهم: «لكل قوم لسان» أي لغة. وقد وصفها ابن جني بأنها «أصوات يعبر بها كل قوم عن أغراضهم»\autocite[33]{ibnjinni}، وهذا التعريف يضع الأصوات في قلب اللغة ويربطها بالتواصل الاجتماعي الهادف، وهو ما يُعدّ اليوم من أسس التعريف اللساني الحديث. وللإشارة، فإن المصطلح الأوروبي \textenglish{Linguistics} يعود بدوره إلى الجذر اللاتيني \textenglish{lingua} أي اللسان، مما يكشف عن توافق لافت بين المعجمين العربي والغربي في هذا التأسيس.

\subsubsection*{ثانياً: التعريف الاصطلاحي}

قبل أن نغوص في دهاليز الجريمة والتحقيق، لا بد أن نضع أقدامنا على أرض صلبة؛ فاللسانيات ليست مجرد «نحو» أو «صرف» كما قد يتبادر للذهن، بل هي ذلك العلم الذي يفكك شفرة اللغة البشرية. وقد عرّفها حليلي بأنها «العلم الذي يدرس اللغة دراسة علمية قائمة على الوصف الدقيق للظواهر اللغوية ومعاينة الوقائع اللغوية كما هي في الاستعمال، بعيداً عن النزعة التعليمية أو إصدار الأحكام المعيارية»\autocite[11]{halili1999}. ويشير وصف اللسانيات بـ«العلم» إلى اعتمادها منهجاً يقوم على أسس موضوعية قابلة للملاحظة والتحقق؛ وهو ما جعلها تتميز تميزاً جوهرياً عن الدراسات النحوية المعيارية القديمة. وتُعرَّف اللسانيات كذلك بأنها العلم الذي يدرس اللغات الطبيعية الإنسانية في ذاتها ولذاتها، سواء أكانت مكتوبة أم منطوقة، مع إعطاء الأولوية للغة المنطوقة باعتبارها المادة الأساسية للتحليل\autocite[19]{hammasa2000}.

بدأت القصة بشكلها الحديث مع فردينان دي سوسير (\textenglish{Ferdinand de Saussure})، الذي أحدث نقلة نوعية في الدراسات اللغوية حين فرّق بين «اللغة» كمنظومة قواعدية مجردة، و«الكلام» كنشاط إنساني ملموس يمارسه الفرد في حياته اليومية (كما هو مبين في الشكل \ref{fig:saussure})\autocite[43]{saussure2008}.

\begin{figure}[H]
\centering
\includegraphics[height=5cm]{Content/02_Chapter1/Fig/Ferdinand de Saussure.png}
\caption{عالم اللسانيات السويسري فردينان دي سوسير، التقطت في أواخر القرن التاسع عشر}
\label{fig:saussure}
\end{figure}

سوسير لم يكتفِ بذلك، بل أحدث ثورة في طريقة تفكيرنا؛ فبدلاً من الانشغال بتتبع تاريخ الكلمات وتطورها عبر العصور، دعا إلى نظرة تزامنية (\textenglish{Synchronic})، أي دراسة اللغة كما هي في لحظتها الراهنة. ثم جاء نعوم تشومسكي (\textenglish{Noam Chomsky}) في الخمسينيات ليضيف بُعداً معرفياً حين ميّز بين «الكفاءة» (\textenglish{Competence})، وهي المعرفة الضمنية بالقواعد، وبين «الأداء» (\textenglish{Performance})، وهو ما يخرج فعلياً من أفواهنا تحت ضغط التعب أو الخوف (انظر الشكل \ref{fig:chomsky}).

\begin{figure}[H]
\centering
\includegraphics[height=5cm]{Content/02_Chapter1/Fig/Noam Chomsky.png}
\caption{عالم اللسانيات الأمريكي نعوم تشومسكي، صاحب نظرية النحو التوليدي وثنائية الكفاءة والأداء}
\label{fig:chomsky}
\end{figure}

وهنا تحديداً، تتقاطع اللسانيات مع عالم التحقيق؛ فالاستجواب ليس مجرد أسئلة وأجوبة، بل هو سياق ضاغط قد يخرج فيه المتحدث عن «كفاءته» المعتادة ليقع في «أداء» يكشف ما يحاول إخفاءه. وهو ما يمنح الخبير الجنائي أدلة لا تقدّر بثمن، إذ أشارت الحافي إلى أن الدراسات اللسانية الحديثة أثبتت أن هذا التناقض بين الكفاءة والأداء هو أحد أهم مداخل كشف الكذب في تحقيقات الشرطة والمحاكم\autocite[4]{alhafi2023}.

ما يميز اللسانيات أنها تخلت عن النظرة المدرسية القديمة التي تحكم على الكلام بـ«الصواب» أو «الخطأ»، واستبدلتها بنظرة وصفية محايدة. هذا التحول هو ما جعلها تخرج من قاعات الدرس لتقتحم مجالات تطبيقية واسعة؛ من تعليم اللغات إلى الترجمة، وصولاً إلى اللسانيات الحاسوبية. ويمكننا أن نتخيل اللسانيات التطبيقية كمظلة كبرى، تندرج تحتها فروع عديدة: اللسانيات التعليمية، وعلم أمراض النطق، واللسانيات الاجتماعية، وبالطبع اللسانيات الجنائية. وما يجمع هذه العائلة هو التركيز على «الاستعمال» الفعلي في سياقات حقيقية بين بشر من لحم ودم\autocite[24]{hammasa2000}.





\subsection{تعريف اللسانيات الجنائية}

\subsubsection*{أولاً: التعريف اللغوي}

نبدأ بتفكيك المصطلح نفسه، لأن ذلك يكشف كثيراً عن ماهية هذا العلم. كلمة «لسانيات» مشتقة من «اللسان»، وهي الترجمة العربية المعتمدة للمصطلح الأوروبي \textenglish{Linguistique} (بالفرنسية) أو \textenglish{Linguistics} (بالإنجليزية)، الذي يعود بدوره إلى الجذر اللاتيني \textenglish{lingua} أي اللسان. وقد وصف العلماء اللسانيات بأنها الدراسة العلمية الموضوعية للغة البشرية من حيث بنيتها ووظائفها ومستوياتها، ونلاحظ هنا كلمة «علمية» التي تفصل اللسانيات المعاصرة عن التأملات الفلسفية القديمة حول اللغة.

أما كلمة «الجنائية» فهي نسبة إلى «الجناية»، والجناية في اللغة العربية مأخوذة من الفعل «جنى يجني» الذي يعني أصلاً قطف الثمار ثم توسّع مجازياً ليدل على اقتراف الذنب. والعلاقة بين المعنيين ليست عشوائية: فمن يقطف ما ليس له يأخذ ما لا يستحق، وتلك هي الجريمة في أبسط تجلياتها. وفي لسان العرب لابن منظور جاء: «جنى عليه جنايةً: أي أذنب ذنباً يُؤاخذ به». فالمصطلح العربي يحمل في بنيته الصرفية فكرة «الاعتداء» و«المؤاخذة»، وهو ما يضعه في قلب المنظومة القانونية\autocite[13]{omar1998}.

في اللغة الإنجليزية، يُستعمل مصطلح \textenglish{Forensic Linguistics}. وكلمة \textenglish{Forensic} تعود في أصلها التاريخي إلى المنتديات الرومانية العامة (\textenglish{Forum})، حيث كانت المناظرات والمحاكمات تُجرى علانية أمام الجمهور. وقد تطور هذا المصطلح تدريجياً ليشمل كل ما له علاقة بالخبرة الفنية أمام القضاء وإجراءات التحقيق المعاصرة، وهو ما يوضح الارتباط الوثيق بين اللغة وسلطة القانون. وقد أشارت الحافي إلى هذه الإشكالية الاصطلاحية مؤكدةً أن تعدد المسميات العربية «أداة لكشف مدى تشابك العلم بالمنظومة القانونية والجنائية في آنٍ واحد»\autocite[4]{alhafi2023}.

ما يستوقفنا هنا — وهي ملاحظة لا نجدها عادة في الكتب المترجمة — أن المصطلح العربي «جنائية» أضيق من المصطلح الإنجليزي \textenglish{Forensic}. فالأول يُحيل مباشرة إلى الجريمة، بينما الثاني يُحيل إلى كل ما يتصل بالقضاء بما فيه النزاعات المدنية والتجارية. وقد تنبّه لهذا الفارق عدد من الباحثين العرب كما سنرى في مسألة المصطلح لاحقاً\autocite[31]{alanati2023}.

\subsubsection*{ثانياً: التعريف الاصطلاحي}

إذا انتقلنا من اللغة إلى الاصطلاح، نجد أن هذا العلم تبلور من خلال ممارسات تطبيقية رائدة. ولعل من أبرز البدايات الموثقة ما قام به جان سفارتفيك عام 1968 في دراسته الشهيرة لقضية تيموثي إيفانز، حيث أثبت بالدليل اللساني أن الاعترافات المنسوبة للمتهم لم تكن من لغته الأصلية، بل أُدخلت عليها تعديلات. وقد أوردت الحافي تفاصيل هذه القضية الرائدة بوصفها نقطة التحول الفعلية التي أعطت اللسانيات الجنائية اعترافاً أكاديمياً ومهنياً بالغ الأثر\autocite[5]{alhafi2023}.

وقد عرّفه جون أولسون بأنه «تطبيق المعرفة اللسانية على البيئة القانونية»، وهو تعريف يتسم بالشمولية ليتسع لكل التقاطعات بين اللغة والقانون\autocite[1]{alhaqbani2008}. بينما ركّز باحثون آخرون مثل كولتهارد وجونسون على الجانب الإجرائي، فاعتبروه «دراسة اللغة بوصفها دليلاً» يُعتدّ به أمام القضاء\autocite[5]{alhafi2023}. هذا التحول في النظر إلى اللغة من مجرد وسيلة تواصل إلى «دليل مادي» هو حجر الزاوية في هذا التخصص.

\begin{figure}[H]
\centering
\includegraphics[height=6.5cm]{Content/02_Chapter1/Fig/Forensic Linguistics.jpg}
\caption{غلاف كتاب «اللسانيات الجنائية» لجون جيبونز (2003)، أحد أهم المراجع التأسيسية في رسم حدود التخصص}
\label{fig:gibbons_book}
\end{figure}

من زاوية أخرى، يشدد ماكمينامين على البعد العلمي والمنهجي لهذا العلم، ويرى أن القيمة الحقيقية تكمن في القدرة على التكرار والاختبار الميداني بعيداً عن الانطباعات الذاتية\autocite[34]{alassimi2018}. ونحن نؤيد هذا الطرح تماماً؛ إذ لا يمكن للقضاء أن يبني أحكاماً على «حدس» لغوي غير منضبط. كما تبرز اللسانيات الوظيفية هنا لتفسير كيف تتأثر بنية اللغة بالوظيفة التي تؤديها في سياق المحاكمة، وهو ما يجعل لكل طرف في العملية القضائية لغة مختارة بعناية تتناسب مع دوره الإجرائي\autocite[107]{halliday2014}.

وعلى الصعيد العربي، لا بد من الإشارة إلى جهود بحثية حديثة بدأت تؤسس لهذا العلم محلياً. فقد أشارت الباحثة طعبة في مقالتها المنشورة بمجلة المحترف للعلوم الإنسانية (2022م) إلى أن اللسانيات الجنائية «لا تزال في مرحلة الطرح تسميةً ودراسةً» في الدول العربية\autocite[60]{taaba2022}. ونبّهت إلى أن المشكلة ليست فقط غياب التنظير بل غياب التواصل بين اللسانيين ورجال القانون في العالم العربي — وهي ملاحظة نتفق معها تماماً.

أما العصيمي فقد قدّم تعريفاً يحاول الجمع بين النظرية والتطبيق، إذ رأى أن اللسانيات الجنائية «فرع من فروع علم اللغة التطبيقي يقوم على دراسة وتحليل البيانات اللغوية المصاحبة لوقوع الجريمة بهدف تحديد هوية الجاني أو المتهم»\autocite[25]{alassimi2018}. والعبادي وسّع النطاق فأضاف بُعداً دولياً حين رأى أن اللسانيات الجنائية تشمل «دراسة النصوص القانونية الدولية ومعاهدات السلام ووثائق المحاكم الجنائية الدولية»\autocite[90]{alabadi2020}، مما يفتح الباب أمام استعمال هذا العلم في جرائم الحرب. بينما حصره الطيب في غاية أضيق وهي «تحديد هوية الجاني أو المتهم»\autocite[267]{altayeb2008}، وهو تعريف ضيّق لكنه يعكس الممارسة الفعلية في أغلب القضايا.

ولا تكتمل الصورة دون الإشارة إلى الحافي التي أكدت في دراستها أن جوهر اللسانيات الجنائية هو «الاحتكام إلى اللغة والاستعانة بها بوصفها وسيلة من وسائل إثبات التهمة أو نفيها»\autocite[5]{alhafi2023}. وهذا التعريف البسيط في ظاهره عميق في مغزاه: اللغة ليست مجرد أداة وصف بل هي ذاتها محلّ الحكم.

\begin{table}[H]
\centering
\caption{مقارنة بين التعريفات الرئيسية للسانيات الجنائية}
\begin{tabular}{r r r r}
\toprule
\textbf{الباحث} & \textbf{التركيز الأساسي} & \textbf{النطاق} & \textbf{الطبيعة} \\
\midrule
سفارتفيك (1968) & التحليل النحوي القضائي & محدود (ممارسة) & تطبيقية \\
أولسون (2008) & البيئة القانونية & واسع & تطبيقية \\
كولتهارد وجونسون (2007) & اللغة كدليل & المحاكم والتحقيقات & عملية \\
ماكمينامين (2002) & الأسلوبية الجنائية & علمي بحت & منهجية \\
العناتي (2023) & اللسانيات التطبيقية & أكاديمي عربي & نظرية + تطبيقية \\
العصيمي (2018) & اللغة والقانون والجريمة & شامل & نظرية + تطبيقية \\
الطيب (2008) & تحديد هوية الجاني & الجريمة فقط & تطبيقية بحتة \\
العبادي (2020) & القانون الجنائي الدولي & دولي & تطبيقية + قانونية \\
\bottomrule
\end{tabular}
\end{table}

ما يتّضح من هذا الاستعراض أن التعريفات تتوزع على طيف واسع: من أضيقها (الطيب) إلى أوسعها (العبادي). والسبب في رأينا ليس ضبابية في التخصص بل تعدد في مداخل النظر إليه — فمن ينطلق من الممارسة المخبرية يُضيّق، ومن ينطلق من الرؤية الأكاديمية يُوسّع. ومن هذا المنطلق، نقترح تعريفاً إجرائياً نعتمده في مذكرتنا: \textbf{اللسانيات الجنائية فرع من اللسانيات التطبيقية يُعنى بتقاطع اللغة مع المنظومة القانونية والقضائية، فيشمل تحليل الأدلة اللغوية لكشف الجرائم وتحديد هوية أصحاب النصوص، وفحص النصوص التشريعية لبيان مقاصدها، وتقييم لغة المحاكم لضمان تواصل عادل بين أطراف العملية القضائية.}

يصبح جلياً أن هذا العلم يقوم على ما يُعرف بـ«النمط اللغوي الفردي» (\textenglish{Idiolect})، وهي القاعدة الذهبية التي تفيد بأن أسلوب كل شخص في ترتيب مفرداته وبناء جمله يشكل نمطاً فريداً يقارب بصمة الإصبع في تفردها. يعمل الخبير على مستويات لغوية متعددة لاستخلاص هذه البصمة، تبدأ من أكثرها مادية وحسية وهو \textbf{المستوى الصوتي} الذي يُحلل طبقات الصوت ونطق الأصوات وما تحمله من آثار لهجية جغرافية. ثم ينتقل الخبير إلى \textbf{المستوى الصرفي والاشتقاقي} حيث تتجلى بصمة الكاتب في طريقة بناء كلماته وميله إلى صيغ اشتقاقية دون غيرها. أما \textbf{المستوى النحوي والتركيبي} فهو من أغنى المستويات إخبارياً، إذ يكشف طول الجمل المعتادة لدى الشخص ومدى تعقيدها أو إيجازها — وكثيراً ما يُفاجأ الخبراء بأن الناس يتمسكون بهذه الأنماط حتى حين يحاولون التخفي. يُكمل ذلك \textbf{المستوى المعجمي} المتعلق بالرصيد اللفظي والمفردات المنتقاة، ثم \textbf{المستوى الدلالي} الذي يفحص المعاني الصريحة والعلاقات الدلالية بين العناصر. ولا يكتمل التحليل دون \textbf{المستوى التداولي} الذي يُعنى بالمقاصد المبطنة والأفعال الكلامية غير المباشرة، وأخيراً \textbf{المستوى الخطابي} الذي يُمثل البنية الكلية للنص ومدى انسجامه وترابطه — وهو ما يُعطي الخبير الصورة الأشمل عن أسلوب صاحب النص، ويلخص الشكل (\ref{fig:diagram_levels}) هذه المستويات السبعة.\autocite[18]{kuwait2023}

\diagramLevels

وتستحق مقولة «البصمة اللغوية» تأملاً أعمق. فخلافاً للبصمة الإصبعية التي ثابتة تماماً ولا تتغير مع الزمن، تتأثر البصمة اللغوية بمتغيرات عديدة: المرحلة العمرية، والسياق الاجتماعي، وحالة الضغط النفسي، والهدف من الكتابة. فالشخص ذاته يكتب بأسلوب مختلف حين يراسل صديقاً أو حين يكتب تقريراً مهنياً. ومن ثمّ فإن عمل الخبير الجنائي الحقيقي يقوم على استخلاص «الثوابت الأسلوبية» التي تظهر في كل هذه السياقات وتبقى رغم التغيرات السطحية. وقد صنف جون جيبونز هذه التقاطعات في مجالات محددة (كما هو واضح في الشكل \ref{fig:diagram_gibbons}).

\diagramGibbons

\subsection{المصطلحات المرتبطة باللسانيات الجنائية}

تحيط باللسانيات الجنائية منظومة من المصطلحات الفرعية والمجاورة التي يُثير الخلطُ بينها وبين المفهوم الأصل إشكاليةً اصطلاحية حقيقية في الأوساط الأكاديمية العربية. وكثيراً ما يستعمل الباحثون المبتدئون في هذا الميدان مصطلح «اللسانيات القانونية» و«اللسانيات الجنائية» بوصفهما مترادفين، في حين أن ثمة حدوداً فاصلة بينهما لا تخفى على المتخصصين.\autocite[12]{doha2023}

\textbf{علم الأسلوب الجنائي (\textenglish{Forensic Stylistics}):} وهو الفرع المتخصص في دراسة الأنماط الكتابية المكتوبة لتحديد ما إذا كان نصان مجهولا المصدر يعودان لكاتب واحد، كما يبرز بوضوح في تحليلات رسائل التهديد أو الوصايا المطعون في صحتها. وقد أشارت الحافي إلى أن الأسلوبية الجنائية «تستعير أدواتها من البلاغة والنقد الأدبي وتُعيد توظيفها في السياق القانوني»، وهو ما يجعلها جسراً حقيقياً بين الأدب والقانون\autocite[7]{alhafi2023}.

\textbf{الصوتيات الجنائية (\textenglish{Forensic Phonetics}):} يُعنى حصراً بالجانب المنطوق للغة. يعتمد على أجهزة التحليل الطيفي (\textenglish{Spectrogram}) لقياس التردد الأساسي، وفترات الصمت، وسرعة الكلام لمطابقة صوت المتهم بتسجيل الجريمة. وهو الأكثر شهرة في التحقيقات الجنائية لكثرة استعمال التسجيلات الهاتفية\autocite[10]{alhafi2023}. ويعمل الخبير هنا كـ«مهندس صوتي» يُحلل الموجات الصوتية بدقة موضوعية، بعيداً عن الانطباعات الشخصية التي كانت سمة التحليل القديم.

\textbf{لسانيات المتون الجنائية (\textenglish{Forensic Corpus Linguistics}):} منهجية تعتمد على قواعد بيانات نصية ضخمة لمعرفة تواتر استعمال كلمة أو عبارة معينة ديموغرافياً، وتفيد في تفسير معاني المصطلحات المستخدمة من قبل شبكات الإجرام\autocite[119]{alassimi2018}. فحين يُدّعى أن عبارة «اعتنِ بنفسك» في رسالة ما هي تهديد مبطّن، يستعمل الخبير متون المراسلات لمعرفة التواتر الإحصائي لهذه العبارة في سياقات التهديد مقابل سياقات الحرص الودّي.

\textbf{تحليل الخطاب القانوني (\textenglish{Legal Discourse Analysis}):} يهتم بطبيعة التواصل المؤسساتي في قاعات المحاكم؛ كالتفاعل بين القاضي والمتهم، والطرق الاستجوابية التي يستخدمها المحامون لتوجيه الإجابات. ودرجة التكافؤ أو اللاتكافؤ بين طرفي الخطاب تمثل محوراً بحثياً خاصاً.

\textbf{علم اللهجات الجنائي (\textenglish{Forensic Dialectology}):} يدرس اللهجات بوصفها أدلة على الأصل الجغرافي أو الاجتماعي للمتكلم، ويُستخدم في قضايا تحديد جنسية اللاجئين (\textenglish{LADO}) وفي تضييق دائرة المشتبه بهم جغرافياً. وقد أشارت الحافي إلى أن هذا الفرع يُمكّن من «إثبات أو نفي صحة حديث مسجل لمتهم وذلك بناء على سمات لهجته الجغرافية والاجتماعية»\autocite[9]{alhafi2023}، وهو ما يجعله أداةً بالغة الأهمية في ملفات اللاجئين القانونية.

إلى جانب هذا التفرع، واجه الباحثون العرب إشكالية في ترجمة المصطلح الأم (\textenglish{Forensic Linguistics}). فبينما يطغى «اللسانيات الجنائية» لكونه يعكس الارتباط المباشر بالجرائم والمحاكم الجنائية، رأى الحقباني مترجم كتاب أولسون أن «القضائي هو الترجمة الصحيحة» لأنه يشمل «اللغة والجريمة والقانون» جميعاً دون الاقتصار على الجانب الجنائي\autocite[1]{alhaqbani2008}. واستعمل آخرون «علم اللغة الجنائي» وهو مجرد اختلاف في تقديم اللسانيات كعلم مستقل أم فرع. واختيارنا لـ«اللسانيات الجنائية» في مذكرتنا نابع من شيوعه من جهة، واعتماده من قبل مؤسسات لغوية رسمية كمجمع الملك سلمان العالمي للغة العربية من جهة أخرى.\autocite[7]{doha2023}

\begin{table}[H]
\centering
\caption{التمييز بين المصطلحات المرتبطة باللسانيات الجنائية}
\begin{tabular}{r r r}
\toprule
\textbf{المصطلح} & \textbf{المجال} & \textbf{العلاقة بالجنائيات} \\
\midrule
اللسانيات الجنائية & التطبيق الجنائي الشامل & المفهوم الأشمل \\
اللسانيات القانونية & اللغة والقانون عموماً & أوسع، يشمل الجنائية \\
التحليل الجنائي للنص & النصوص المكتوبة فقط & مجال فرعي \\
الصوتيات الجنائية & الأدلة الصوتية & مجال فرعي مستقل \\
لغة القانون & النصوص التشريعية & مجاور لكن مختلف \\
علم الأسلوب الجنائي & الخصائص الكتابية الفردية & مجال فرعي تطبيقي \\
علم اللهجات الجنائي & الأصل الجغرافي للمتكلم & مجال فرعي ظرفي \\
\bottomrule
\end{tabular}
\end{table}



%% ======================================================
\section{علاقة اللسانيات الجنائية بالعلوم الأخرى}
%% ======================================================

تندرج اللسانيات الجنائية بامتياز ضمن الدراسات البينية أو العابرة للتخصصات (\textenglish{Interdisciplinary Studies}). فهي لم تنشأ في فراغ، بل هي ثمرة تزاوج معرفي بين علوم تبدو متباعدة ظاهرياً، لكنها تتحد في هدف أسمى هو كشف الحقيقة وإرساء العدالة. وهذه الطبيعة البينية ليست نقيصة — كما قد يظنها البعض — بل هي من أكبر مزاياها؛ إذ تمكّنها من استعارة المنهجيات الأكثر ملاءمة لكل سياق وتوظيفها بمرونة لا يتيحها أي علم ضيق التخصص.\autocite[4]{cerist2023}

\subsection{علاقتها بالقانون}

تمثل صلة اللسانيات بالعلوم القانونية العمود الفقري لهذا التخصص. القانون في طبيعته هو «مهنة الكلمات»؛ المشرّع يصوغ القوانين بالكلمات، والقاضي يفسرها ويصدر حكمه بها، والمحامي يرافع معتمداً على بلاغتها. فكل فعل قانوني يمر بالضرورة عبر قناة لغوية\autocite[85]{alabadi2020}.

تكتسي اللسانيات أهميتها هنا من حيث أنها تضبط المفاهيم القانونية وتحلل دلالاتها. فغالباً ما تنشأ النزاعات بسبب الغموض الدلالي أو التركيبي في صياغة العقود أو التشريعات. ولنأخذ مثالاً واقعياً: كلمة «أو» في لغة التشريع تُستعمل أحياناً بمعنى الجمع الإضافي «و»، وأحياناً بمعنى الاستئناف الحصري (إما هذا أو ذاك)، مما يُفضي إلى خلافات تفسيرية حقيقية. يتدخل اللساني لتقديم قراءة موضوعية ترتكز على قواعد النحو والدلالة لترجيح المعنى المقصود.\autocite[20]{nader2023}

كما برز بوضوح دور اللسانيات من خلال «حركة تبسيط اللغة القانونية» (\textenglish{Plain Language Movement})\autocite[98]{alabadi2020}، التي تسعى لفك طلاسم لغة التشريعات المعقدة لضمان حق المواطن البسيط في فهم ماله وما عليه. وقد أثبتت الدراسات المقارنة أن كثيراً من الأحكام الجنائية جاءت مجحفة لأن المتهم لم يفهم نصاً قانونياً كان يستلزم معرفته. والمنطق وراء ذلك بسيط وعميق في الآن ذاته: كيف يمكن محاسبة شخص على خرق قانون لم يستطع فهمه أصلاً بسبب تعقيد صياغته؟ إن العدالة التي لا تُفهم لغتها عدالة منقوصة من جوهرها.

ونشير هنا إلى مسألة تفسير العقود التي تُشكّل ميداناً بارزاً لتطبيق اللسانيات الجنائية في القضايا المدنية. فحين يتنازع طرفان على تفسير بند في عقد تجاري، تُقدّم اللسانيات تحليلاً موضوعياً للدلالات المحتملة للنص، آخذةً في الاعتبار المعايير التداولية للتواصل القانوني في السياق والزمان والمكان المعنيين. وقد أصبح هذا التطبيق شائعاً في المحاكم الأمريكية والبريطانية حيث يُستدعى اللسانيون كثيراً في نزاعات العلامات التجارية لتحديد ما إذا كانت علامتان بينهما تشابه مُضلِّل\autocite[24]{alabadi2020}.

\subsection{علاقتها بعلم الاجتماع}

لا يتحدث الإنسان بمعزل عن مجتمعه؛ فطبقته الاجتماعية، ومستواه التعليمي، وحتى بيئته الجغرافية تُطبع على لسانه. وهنا تتقاطع اللسانيات الجنائية مع علم الاجتماع اللغوي (\textenglish{Sociolinguistics}) بطريقة حيوية في المحاكم.

من أبرز انعكاسات هذه العلاقة ما واجهته المحاكم الأسترالية حول الصعوبات الجمة التي يعاني منها السكان الأصليون في فهم لغة المحققين لتباين الخلفيات الاجتماعية والثقافية\autocite[92]{alabadi2020}. يقرّ المتهم أحياناً بذنبه لا عن اقتناع، بل مسايرة للمحققين لأسباب ثقافية تعتبر معارضة السلطة عيباً اجتماعياً. تفكك اللسانيات الجنائية هذه الشيفرة وتثبت للقضاء أن الإجابة بـ«نعم» قد تعني في الثقافة الأصلية للمتهم مجرد إقرار بالسماع وليس اعترافاً بالفعل، مما يحمي الفئات المجتمعية الهشة من إجحاف الأحكام القائمة على سوء الفهم الثقافي.

وثمة وجه آخر لهذه العلاقة يتجلى في تحديد الأصل الاجتماعي من خلال الكلام. فعلماء الاجتماع اللغوي يعلمون جيداً أن أنماط الكلام تختلف من طبقة اجتماعية إلى أخرى — في اختيار المفردات، وفي بنية الجمل، وفي التوقف والتردد. وهذا المعرفة تُستثمر في اللسانيات الجنائية حين يُحاول الخبير تحديد ملامح الكاتب أو المتكلم المجهول. فرسالة الفدية التي تُوحي بالتعليم العالي والكتابة الرسمية وأخرى تفيض بمعجم العامة ولغة الشارع — كلتاهما تحمل تفاصيل اجتماعية لا يستطيع تفسيرها إلا من يُحسن قراءة اللغة كظاهرة اجتماعية.

وتزداد أهمية الحساسية الاجتماعية في سياق التهديدات الإلكترونية والتواصل الرقمي؛ إذ باتت الجرائم اللغوية لا تُفهم إلا في سياقها الاجتماعي الدقيق. فعبارة تُكتب بحروف كبيرة (\textenglish{ALL CAPS}) في رسالة إلكترونية لها دلالة اجتماعية واضحة لمن يفهم أعراف التواصل الرقمي، وهي دلالة لا تُستقرأ إلا من خلال العدسة الاجتماعية اللغوية.

\subsection{علاقتها بعلم النفس الجنائي}

ترتبط لغة الفرد ارتباطاً وثيقاً بحالته النفسية والانفعالية. حين يُجري خبراء النفس تحليلاً لشخصية المجرم وميولاته، فإن المادة الأولية التي يدرسونها هي ما يصرح به قولاً أو كتابة. يتجلى تقاطع اللسانيات بعلم النفس (\textenglish{Psycholinguistics}) الجنائي بقوة في عمليات استجواب المتهمين، حيث يسعى المحقق لاكتشاف التناقض المعرفي والانفعالي في كلام المشتبه به لكشف الخداع والكذب.

وقد طوّر علماء النفس اللغوي نماذج نظرية تُفسّر آليات الكذب في اللغة. من أبرزها نموذج «إدارة الانطباع» (\textenglish{Impression Management}) الذي يقول إن الكاذب يبذل جهداً ذهنياً إضافياً لتحضير روايته، مما يُفرز في كلامه مؤشرات يمكن رصدها: جمل أقصر، تردد أكثر في المواضع الحساسة، انتقال مفاجئ من الضمير «أنا» إلى «الناس» أو «المرء» حين يُدلي بمعلومات افترى فيها. هذه المؤشرات يجمع بين التحليل النفسي واللساني معاً.

كما يظهر هذا التعاون جلياً في تحليل الرسائل الانتحارية؛ فالخبير اللساني يبحث في النص عن مؤشرات صدق العاطفة والانكسار النفسي الذي يميز المقدمين على الانتحار الفعلي. وقد أشارت طعبة في دراستها إلى هذا الجانب التطبيقي المهم، موضحةً أن الخبير اللساني يستطيع التمييز بين النص الصادق ونص مَن يُقدّم حالة انتحارية مزوّرة لإخفاء جريمة، من خلال المؤشرات الأسلوبية كالطول ودرجة التفاصيل والمعجم المستخدم\autocite[56]{taaba2022}.

وفي مجال الشهادات أمام المحاكم، تُفيد الأبحاث النفسية اللغوية في فهم كيفية تشكّل ذاكرة الشهود وقابليتها للتأثر بالتوجيه اللغوي للأسئلة. فالسؤال «بأي سرعة كانت السيارة تسير حين اصطدمت؟» يُنتج إجابات تُقدّر السرعة بأعلى مما يُنتجه السؤال «بأي سرعة كانت السيارة تسير حين مرّت أمامك؟» رغم أنهما يصفان الحادثة ذاتها. هذا التأثير، الذي وثّقته أبحاث علم النفس المعرفي الرائدة، يُعدّ من أهم الإسهامات في نقطة تقاطع علم النفس واللسانيات الجنائية.

\subsection{علاقتها بعلم الأصوات}

إذا كانت البصمة الإصبعية الدعامة الأولى للأدلة الجنائية البيولوجية، فإن علم الأصوات (\textenglish{Phonetics}) هو المعادل المادي والفيزيائي للغة في قاعات المحاكم. لا تتعامل اللسانيات الجنائية مع الصوت بانطباعية، بل تستعير من الفونيتيكا أدواتها المخبرية الدقيقة لعزل المويجات الصوتية وفحص تردداتها وشدتها.

تساعد هذه الشراكة الأجهزة الأمنية في تحليل شرائط التهديد وتضييق دائرة المشتبه بهم بناءً على المخرج الصوتي (كالهمزة والقاف)، أو سرعة النطق النمطية، أو حتى وجود اختلالات تنفسية أثناء الكلام تُشير لعوارض صحية محددة أو حالة توتر شديد. وقد أسهمت علاقة العلمين في تطوير برمجيات قوية قادرة اليوم على تصفية التسجيلات المشوشة لتفريغ المحادثات بدقة\autocite[12]{alhafi2023}.

ولا بد من الإشارة إلى الجدل العلمي المتعلق بمصطلح «البصمة الصوتية» (\textenglish{Voiceprint}). فبينما يرى بعض الخبراء أن الصوت الإنساني فريد كالبصمة الإصبعية تماماً، يُحاجج آخرون بأن الصوت متغير بطبيعته: يتأثر بالمرض والتعب والعمر والانفعال، مما يجعل الحديث عن «بصمة ثابتة» يحتاج إلى تحفظات علمية جوهرية. وقد لجأ علماء الصوتيات الجنائية إلى إطار أكثر حذراً يتحدث عن «احتمال التطابق» بدلاً من «تطابق حاسم»، تماشياً مع متطلبات الشهادة العلمية أمام القضاء. وقد أشارت طعبة إلى أن «الأصوات قد تتشابه، لكنها كالبصمة لا تتطابق أبداً»، مؤكدةً ضرورة اعتماد معيار احتمالي لا قطعي في الخبرة الصوتية الجنائية\autocite[57]{taaba2022}.

ومن أحدث التطبيقات في هذا المجال استخدام تقنيات التعلم العميق (\textenglish{Deep Learning}) في تحليل الصوت. فأصبح بالإمكان اليوم التعرف على متكلم من عينة صوتية قصيرة جداً، وتمييز الصوت المُعدَّل أو المُقنَّع رقمياً، بل وكشف التزوير الصوتي المُنتج بتقنية \textenglish{Deepfake}. وهذا الأخير يطرح تحديات أخلاقية وقانونية بالغة الخطورة؛ إذ أصبح بالإمكان توليد تسجيلات صوتية وهمية مقنعة جداً تُنسب لأشخاص حقيقيين. ومما يستوجب العناية الفائقة أن تواكب اللسانيات الجنائية هذا التطور وتُطور أدوات موثوقة لكشف هذه التزويرات قبل أن تصبح سلاحاً في يد المتهمين لتلفيق أدلة صوتية مزورة.



%% ======================================================
\section{أهمية الدليل اللغوي في السياق الجنائي}
%% ======================================================

\subsection{اللغة كأداة لارتكاب الجريمة}

لا تقتصر الجريمة على الفعل الجسدي الملموس؛ فاللغة قادرة وحدها على إحداث ضرر قانوني بليغ. ما يضفي أهمية بالغة على اللسانيات الجنائية هو قدرتها على تفكيك هذه الجرائم اللغوية. يصنف الباحثون جرائم لا تقع إلا باللغة، مثل السب والشتم والقذف والتشهير والاحتيال والتحريض. فقد ميّزت الحافي بين جرائم لغوية مباشرة «أداتها اللغة، حيث يرتكب الجاني الجريمة بواسطة اللغة إما منطوقة أو مكتوبة» -- وتشمل التهديد والرشوة وشهادة الزور -- وجرائم لغوية غير مباشرة «لم تكن اللغة فيها برمتها وإنما اللغة عامل مهم فيها»\autocite[14]{alhafi2023}. كما شكّل التهديد وطلبات الفدية حقلاً خصباً للتحليل الأسلوبي لكشف هوية المرسل\autocite[54]{taaba2022}.

الأعقد من ذلك هي «الرشوة اللغوية» والتحريض. فالراشي نادراً ما يستخدم عبارة مباشرة فجة، بل يلجأ للمواربة والتداولية (كأن يضع مالاً ويقول «هذا لتشرب القهوة»). القاضي يحتاج لخبير محايد يفكك هذا التضمين لإثبات القصد الجنائي المختبئ خلف قناع اللغة المهذبة. وقد أسس روجر شوي لهذا المنهج التطبيقي مبكراً في المحاكم، حين فضح أساليب الضغط اللغوي والاستدراج الذي كانت تمارسه الجهات الأمنية لتوريط المتهمين في اعترافات ملفقة\autocite[34]{alassimi2018}.

ومن الجرائم اللغوية التي تستوجب تحليلاً متخصصاً جريمة «انتحال الشخصية» (\textenglish{Impersonation}) في الفضاء الرقمي؛ إذ يُحاكي المجرم أسلوب ضحيته الكتابي لخداع الآخرين. والكشف عن هذا الانتحال يقوم أساساً على مقارنة البصمة اللغوية للنصوص الحقيقية للضحية مع النصوص المنتحلة، بحثاً عن الانزياحات الأسلوبية التي تكشف أن مؤلفاً مختلفاً وراء القلم.

وتبرز جريمة «التدليس اللغوي» (\textenglish{Linguistic Fraud}) في عالم الأعمال حين تُصاغ عقود أو تقارير مالية بعبارات ملتبسة مقصودة تُخفي حقائق مضرة بالطرف الآخر. وقد استُعين باللسانيين في قضايا الاحتيال المالي الكبرى للكشف عن هذا الغموض المتعمد وإثباته أمام القضاء.

\subsection{اللغة كدليل مادي أو قرينة}

من المفيد التوقف عند الفارق القانوني بين الدليل والقرينة في التعامل مع المعطيات اللغوية. الدليل هو الإثبات القاطع الذي يحسم النزاع، أما القرينة فهي مؤشر قوي يحتاج لما يسنده.

في الجرائم اللغوية الخالصة (كالتهديد بالقتل المبرق عبر رسالة)، الرسالة نفسها هي مسرح الجريمة وأداة الجريمة والدليل القاطع عليها معاً. في المقابل، حين تتعلق القضية بجريمة قتل، ويُعثر على مذكرات للمتهم يهاجم فيها الضحية، فإن اللغة هنا ترتقي لمرتبة «القرينة الدامغة» التي تساعد في إثبات دافع الكراهية، لكنها لا تثبت منفردة لحظة القتل المادية\autocite[22]{alassimi2018}.

وفي هذا السياق، تُفرّق المراجع العلمية بين ثلاثة مستويات للدليل اللغوي: الدليل المباشر (حين تكون الجريمة ذاتها لغوية كالتشهير)، والدليل الظرفي (حين تُستعمل اللغة لتحديد الفاعل أو الدافع)، والدليل الإسنادي (حين تُستعمل اللغة لربط شخص بعينه بعمل إجرامي). وقد أشارت الحافي إلى أن إغفال هذه التمييزات يُفضي إلى تقديم الدليل اللغوي بصورة مُضخّمة أو مُهوَّنة أمام القضاء\autocite[13]{alhafi2023}. ويمكن تتبع هذه الطبيعة المزدوجة بوضوح عبر المخطط البياني (الشكل \ref{fig:diagram_evidence}).

\diagramEvidence

\subsection{حجية الدليل اللغوي أمام القضاء}

يطرح الاعتماد المتزايد على الخبرة اللسانية سؤالاً مفصلياً حول مدى مشروعية وحجية هذا التقرير أمام قاضي الحكم. لا يقبل القضاء الاستنتاجات المبنية على الحدس أو المشاعر تحت مسمى «الخبرة». ولهذا، يلتزم اللسانيون الجنائيون في الدول المتقدمة بمعايير صارمة، أشهرها «معيار دوبرت» (\textenglish{Daubert Standard}) الذي أقرته المحاكم الأمريكية\autocite[96]{alabadi2020}.

يقضي هذا المعيار بأن الشهادة العلمية لا تُقبل ما لم تكن خاضعة لمنهجية قابلة للقياس والاختبار، ومنشورة في دوريات محكّمة، وذات نسبة هامش خطأ معروف. بعبارة أخرى، على اللساني ألا يقول «يبدو لي أن فلان هو الكاتب»، بل عليه أن يصرح: «تطابق الأسلوب بين رسالة التهديد والنصوص المرجعية للمتهم بلغ نسبة إحصائية تدل بقوة على كونه المؤلف الأصلي»\autocite[57]{alassimi2018}.

أما في النظام القانوني البريطاني فيُعمل بما يُعرف بـ«معيار دوبرت المعدَّل» أو «اختبار المحكمة العليا» الذي يمنح القاضي صلاحية أوسع في تقييم موثوقية الخبرة العلمية. وقد أفضى هذا المرونة النسبية إلى قبول شهادات لسانية كانت ستُرفض في المحاكم الأمريكية الأكثر صرامة. والنقاش حول المعيار الأنسب لا يزال دائراً في الأوساط الأكاديمية والقضائية الغربية.

وثمة بُعد أخلاقي لا يمكن إغفاله في سياق حجية الدليل اللغوي. فالخبير اللساني وإن كان يعمل بأدوات علمية موضوعية، إلا أنه بشر قد تتأثر استنتاجاته بالجهة التي استأجرته وبالمعلومات المسبقة التي تلقّاها عن المتهم. ولهذا يُطالب كثيرون من المختصين بأن يُكشف لكل خبير في المحكمة عن الجهة الممولة لشهادته، وأن يُتاح للطرف الآخر إمكانية الطعن في منهجيته بالاستعانة بخبير منافس.

ويجب التأكيد ختاماً على ضرورة التزام الخبير بحدود تخصصه؛ اللساني الجنائي ليس طرفاً متحيزاً ولا يصدر أحكاماً ببراءة المتهم أو تجريمه. واجبه يقتصر على تنوير المحكمة بمعطيات علمية دقيقة حول هوية صاحب النص أو حقيقة مقاصده، تاركاً السلطة التقديرية للقاضي الذي يتخذ التدابير النهائية بناء على منظومة الأدلة المتكاملة.



%% ======================================================
\section{مراحل عمل اللساني الجنائي}
%% ======================================================

\subsection{من الجريمة إلى الحكم: المسار الإجرائي}

يخضع عمل اللساني الجنائي لمسار إجرائي صارم يبدأ من لحظة تلقّي الأدلة وينتهي بتقديم شهادة الخبرة أمام القضاء. وفهم هذا المسار ضروري للقارئ كي يُدرك أن الخبرة اللسانية ليست مجرد «قراءة» لنص، بل هي عملية منهجية متعددة المراحل تخضع لضوابط صارمة (يلخص الشكل \ref{fig:diagram_workflow} هذه المراحل المتسلسلة).

\diagramWorkflow

تبدأ العملية بما يمكن تسميته مرحلة \textbf{الاستلام والتوثيق}، وهي أكثر المراحل رتابة وأشدها أهمية في الوقت ذاته. يتسلّم الخبير الأدلة اللغوية — نصوصاً أو تسجيلات أو وثائق — ويوثّق كل ما يحيط بها من معلومات سياقية: متى وُجدت، وأين، وبأي ظروف، ومن تولّى حفظها. هذا التوثيق الممل ظاهرياً ضروري لأن أي ثغرة في سلسلة حفظ الأدلة قد يستغلها الدفاع لاحقاً للطعن في صحتها.

بعد ذلك يتجه الخبير إلى \textbf{التحليل الأولي} الذي يُعطيه صورة سريعة عن طبيعة النص: اللغة المستعملة فصحى أم عامية أم مختلطة، والمستوى التعليمي المُستشف منها، وأي ظواهر لغوية لافتة تستدعي تعمقاً أكبر.

ثم تأتي المرحلة الأكثر حساسية وهي \textbf{جمع النصوص المرجعية} من المشتبه بهم: رسائل سابقة، وثائق شخصية، تسجيلات معروفة الهوية. هذه المرجعيات هي المحك الذي يُقاس إليه النص المجهول، وكلما كانت أوفر وأنوع كان التحليل أكثر موثوقية.

وبعد اكتمال المرجعيات يبدأ \textbf{التحليل المعمّق والمقارن} — وهو القلب النابض للعملية كلها — حيث تُطبَّق الأدوات الكمية والنوعية معاً: قياس الثروة المعجمية، وتحليل توزيع أطوال الجمل، ورصد الأنماط التركيبية المتكررة، ومقارنة التواترات اللغوية. وتختتم العملية بمرحلة \textbf{كتابة التقرير} حيث يُحوّل الخبير نتائجه التقنية إلى صياغة مفهومة للقاضي وهيئة المحلفين، مع التصريح بهامش اليقين والخطأ المحتمل — لأن الخبير الذي يدّعي اليقين المطلق يُشير بذاته إلى ضعف منهجيته.

\subsection{معايير الكفاءة في الخبرة اللسانية الجنائية}

لا يكفي أن يكون الشخص عالماً في اللسانيات لكي يكون خبيراً جنائياً موثوقاً. وهذه نقطة كثيراً ما يُستهان بها. فهذا التخصص يستلزم تقاطع كفاءات متعددة يُعسر الجمع بينها، وإن كان بالإمكان بناؤها بالتدريج.

تبدأ هذه الكفاءات بالجانب التحليلي البحت: القدرة على تطبيق الأدوات اللسانية بمنهجية دقيقة على أنواع النصوص المختلفة، وهو ما يستلزم تكويناً متيناً في لسانيات المتون والصوتيات وتحليل الخطاب. لكن هذه الكفاءة وحدها غير كافية ما لم تصحبها معرفة وافية بالمنظومة القانونية ومتطلباتها الإجرائية. الخبير الذي لا يعرف ما هو «معيار دوبرت» أو ما الذي تعنيه «درجة الإثبات» قانونياً، ستُرفض شهادته لمجرد أنها لا تستوفي الشكل الإجرائي المطلوب، بغض النظر عن متانتها العلمية.

والشيء الذي لا تُدرسه الجامعات في الغالب هو الكفاءة التواصلية: القدرة على ترجمة نتائج التحليل التقني المعقد إلى لغة واضحة مفهومة لقاضٍ أو لهيئة محلفين بلا خلفية لسانية. وقد أكد الحقباني في مقدمته لترجمة كتاب أولسون على هذه النقطة بإلحاح\autocite[1]{alhaqbani2008}، مشدداً على أن الخبير الذي يلجأ إلى مصطلحات تقنية متراكمة أمام القاضي يضر قضيته أكثر مما ينفعها. والمحاكم مليئة بأمثلة على تقارير لسانية متقنة أُهدرت لأن صاحبها عجز عن إيصالها.

وتأتي في الأخير — وأقول في الأخير لصعوبتها لا لأهميتها — الكفاءة الأخلاقية. الخبير ملزم بتقديم الحقيقة التي يرى كاملة حتى حين تكون ضد مصلحة الجهة التي استأجرته. والانحياز المؤسسي — سواء كان واعياً أو لا واعياً — أفسد في تاريخ القضاء شهادات عديدة وجرّ أحكاماً ظالمة. ولعل هذا هو ما جعل كثيراً من الخبراء الغربيين يُؤثرون العمل لصالح المحاكم بوصفهم خبراء محايدين معيّنين قضائياً، بدلاً من العمل لصالح طرف بعينه.



% خلاصة الفصل الأول
\vspace{0.5cm}
\noindent\textbf{خلاصة الفصل الأول:}

يتضح جلياً مما سبق أن اللسانيات الجنائية لم تعد ترفاً أكاديمياً، بل ضرورة إجرائية تفرضها طبيعة الجرائم المعاصرة. انطلقنا من ضبط مصطلحات العلم وتتبع تشعبه في منظومة بينية ثرية تجمع اللغة بالقانون وعلم الاجتماع وعلم النفس وعلم الأصوات. ورأينا كيف أن الدليل اللغوي يأخذ أشكالاً متعددة ويخضع لمستويات متباينة من الحجية والقوة الإثباتية. ثم توقفنا عند المسار الإجرائي للعمل الجنائي اللساني ومعايير الكفاءة التي تحكمه، إدراكاً منا أن الجودة في هذا التخصص ليست مجرد معرفة نظرية بل منهجية تطبيقية واحترام أخلاقي لأبعاد العدالة. وسيقودنا هذا الفهم للإطار النظري إلى الغوص عميقاً في التطورات التاريخية لهذا العلم وتأصيل جذوره في منظومة تراثنا العربي، وهو ما سنفصله في الفصل الثاني.

